% =====================================================================
% ACT II DATA PACK — Recovery Is Parameter Reset, Not Restoration (T3-S5 · 2026-10-19)
% Sources:
%   * R1–R6 definitions: SMPDL_Memo_AuditStructure_and_Removal.md §4.2
%   * Operator lens trio (Restructuring / Part2 OntologyDynamics / OutOfSample)
%   * Instance mapping: corpus rows below cite group IDs (H#) only —
%     every ID traces to its adjudicated report.
% Loaded by main.tex §5 via % =====================================================================
% ACT II DATA PACK — Recovery Is Parameter Reset, Not Restoration (T3-S5 · 2026-10-19)
% Sources:
%   * R1–R6 definitions: SMPDL_Memo_AuditStructure_and_Removal.md §4.2
%   * Operator lens trio (Restructuring / Part2 OntologyDynamics / OutOfSample)
%   * Instance mapping: corpus rows below cite group IDs (H#) only —
%     every ID traces to its adjudicated report.
% Loaded by main.tex §5 via % =====================================================================
% ACT II DATA PACK — Recovery Is Parameter Reset, Not Restoration (T3-S5 · 2026-10-19)
% Sources:
%   * R1–R6 definitions: SMPDL_Memo_AuditStructure_and_Removal.md §4.2
%   * Operator lens trio (Restructuring / Part2 OntologyDynamics / OutOfSample)
%   * Instance mapping: corpus rows below cite group IDs (H#) only —
%     every ID traces to its adjudicated report.
% Loaded by main.tex §5 via % =====================================================================
% ACT II DATA PACK — Recovery Is Parameter Reset, Not Restoration (T3-S5 · 2026-10-19)
% Sources:
%   * R1–R6 definitions: SMPDL_Memo_AuditStructure_and_Removal.md §4.2
%   * Operator lens trio (Restructuring / Part2 OntologyDynamics / OutOfSample)
%   * Instance mapping: corpus rows below cite group IDs (H#) only —
%     every ID traces to its adjudicated report.
% Loaded by main.tex §5 via \input{act2_data}
% =====================================================================

\subsection{Parameter-Bundle Ontology and Six Reset Patterns}\label{sec:act2-rpatterns}

\textbf{Definition.} A law is an intensional conditional parameter
bundle; removal disassembles the bundle; recovery is one of six
parameter re-wiring patterns:

\begin{table}[h]
\centering
\begin{tabular}{llp{0.46\textwidth}}
\toprule
 & \textbf{Pattern} & \textbf{Corpus instances (group IDs)} \\
\midrule
R1 & Same-parameter weakening & $C^2$ regularity $\to$ weak
derivative / Sobolev $H^1$ setting (H18); induction full $\to$
fragment ladder (Big-Five side, H21 corpus neighbor). \\
R2 & Different-parameter substitution & Pasch-ordered $\to$
betweenness/coordinatized routes; Riemann-integral limit exchange
$\to$ monotone-class arguments (H17 side). \\
R3 & Parameter coordinatization & Parallel postulate $\to$ curvature
parameter $K$ (H1/H16); non-Archimedean valuation rank (H34 side). \\
R4 & Parameter pricing & Integral order-exchange three-tier pricing
(Tonelli free zone / Fubini conditional / failure) (H30). \\
R5 & Parameter measurization & Selmer group, Tate--Shafarevich group
as computable obstructions (H46); Brauer--Manin obstruction. \\
R6 & Parameter programmatization & Cut-elimination as algorithm
(H22 side); Knuth--Bendix completion (H31 meta-layer emergence). \\
\bottomrule
\end{tabular}
\caption{Six recovery patterns with corpus anchors. R1/R2 rows carry
fewer adjudicated exemplars than R3--R6 in the current files; the
mapping is stated at the granularity the evidence supports (per
adjudication discipline, no inflation).}
\end{table}

\subsection{Proposition: Reset, Not Restoration}\label{sec:act2-proposition}

\begin{proposition}[Reset-not-restoration; informal wording, final at S5 fill]
In no corpus instance is the recovery product a copy of the removed
law; in every instance the death position of the old law is occupied
by a structure that keeps its books differently (a new concept
re-wired to the disassembled parameters). Typological induction over
all 29 removal groups.
\end{proposition}

\noindent\emph{Evidence anchor.} The 29-group summary table (Act III
data pack, Table~\ref{tab:act3-summary}) lists the emerged concept
for every group; none restorates the removed law; each realizes one
of R1--R6.

% =====================================================================

\subsection{Parameter-Bundle Ontology and Six Reset Patterns}\label{sec:act2-rpatterns}

\textbf{Definition.} A law is an intensional conditional parameter
bundle; removal disassembles the bundle; recovery is one of six
parameter re-wiring patterns:

\begin{table}[h]
\centering
\begin{tabular}{llp{0.46\textwidth}}
\toprule
 & \textbf{Pattern} & \textbf{Corpus instances (group IDs)} \\
\midrule
R1 & Same-parameter weakening & $C^2$ regularity $\to$ weak
derivative / Sobolev $H^1$ setting (H18); induction full $\to$
fragment ladder (Big-Five side, H21 corpus neighbor). \\
R2 & Different-parameter substitution & Pasch-ordered $\to$
betweenness/coordinatized routes; Riemann-integral limit exchange
$\to$ monotone-class arguments (H17 side). \\
R3 & Parameter coordinatization & Parallel postulate $\to$ curvature
parameter $K$ (H1/H16); non-Archimedean valuation rank (H34 side). \\
R4 & Parameter pricing & Integral order-exchange three-tier pricing
(Tonelli free zone / Fubini conditional / failure) (H30). \\
R5 & Parameter measurization & Selmer group, Tate--Shafarevich group
as computable obstructions (H46); Brauer--Manin obstruction. \\
R6 & Parameter programmatization & Cut-elimination as algorithm
(H22 side); Knuth--Bendix completion (H31 meta-layer emergence). \\
\bottomrule
\end{tabular}
\caption{Six recovery patterns with corpus anchors. R1/R2 rows carry
fewer adjudicated exemplars than R3--R6 in the current files; the
mapping is stated at the granularity the evidence supports (per
adjudication discipline, no inflation).}
\end{table}

\subsection{Proposition: Reset, Not Restoration}\label{sec:act2-proposition}

\begin{proposition}[Reset-not-restoration; informal wording, final at S5 fill]
In no corpus instance is the recovery product a copy of the removed
law; in every instance the death position of the old law is occupied
by a structure that keeps its books differently (a new concept
re-wired to the disassembled parameters). Typological induction over
all 29 removal groups.
\end{proposition}

\noindent\emph{Evidence anchor.} The 29-group summary table (Act III
data pack, Table~\ref{tab:act3-summary}) lists the emerged concept
for every group; none restorates the removed law; each realizes one
of R1--R6.

% =====================================================================

\subsection{Parameter-Bundle Ontology and Six Reset Patterns}\label{sec:act2-rpatterns}

\textbf{Definition.} A law is an intensional conditional parameter
bundle; removal disassembles the bundle; recovery is one of six
parameter re-wiring patterns:

\begin{table}[h]
\centering
\begin{tabular}{llp{0.46\textwidth}}
\toprule
 & \textbf{Pattern} & \textbf{Corpus instances (group IDs)} \\
\midrule
R1 & Same-parameter weakening & $C^2$ regularity $\to$ weak
derivative / Sobolev $H^1$ setting (H18); induction full $\to$
fragment ladder (Big-Five side, H21 corpus neighbor). \\
R2 & Different-parameter substitution & Pasch-ordered $\to$
betweenness/coordinatized routes; Riemann-integral limit exchange
$\to$ monotone-class arguments (H17 side). \\
R3 & Parameter coordinatization & Parallel postulate $\to$ curvature
parameter $K$ (H1/H16); non-Archimedean valuation rank (H34 side). \\
R4 & Parameter pricing & Integral order-exchange three-tier pricing
(Tonelli free zone / Fubini conditional / failure) (H30). \\
R5 & Parameter measurization & Selmer group, Tate--Shafarevich group
as computable obstructions (H46); Brauer--Manin obstruction. \\
R6 & Parameter programmatization & Cut-elimination as algorithm
(H22 side); Knuth--Bendix completion (H31 meta-layer emergence). \\
\bottomrule
\end{tabular}
\caption{Six recovery patterns with corpus anchors. R1/R2 rows carry
fewer adjudicated exemplars than R3--R6 in the current files; the
mapping is stated at the granularity the evidence supports (per
adjudication discipline, no inflation).}
\end{table}

\subsection{Proposition: Reset, Not Restoration}\label{sec:act2-proposition}

\begin{proposition}[Reset-not-restoration; informal wording, final at S5 fill]
In no corpus instance is the recovery product a copy of the removed
law; in every instance the death position of the old law is occupied
by a structure that keeps its books differently (a new concept
re-wired to the disassembled parameters). Typological induction over
all 29 removal groups.
\end{proposition}

\noindent\emph{Evidence anchor.} The 29-group summary table (Act III
data pack, Table~\ref{tab:act3-summary}) lists the emerged concept
for every group; none restorates the removed law; each realizes one
of R1--R6.

% =====================================================================

\subsection{Parameter-Bundle Ontology and Six Reset Patterns}\label{sec:act2-rpatterns}

\textbf{Definition.} A law is an intensional conditional parameter
bundle; removal disassembles the bundle; recovery is one of six
parameter re-wiring patterns:

\begin{table}[h]
\centering
\begin{tabular}{llp{0.46\textwidth}}
\toprule
 & \textbf{Pattern} & \textbf{Corpus instances (group IDs)} \\
\midrule
R1 & Same-parameter weakening & $C^2$ regularity $\to$ weak
derivative / Sobolev $H^1$ setting (H18); induction full $\to$
fragment ladder (Big-Five side, H21 corpus neighbor). \\
R2 & Different-parameter substitution & Pasch-ordered $\to$
betweenness/coordinatized routes; Riemann-integral limit exchange
$\to$ monotone-class arguments (H17 side). \\
R3 & Parameter coordinatization & Parallel postulate $\to$ curvature
parameter $K$ (H1/H16); non-Archimedean valuation rank (H34 side). \\
R4 & Parameter pricing & Integral order-exchange three-tier pricing
(Tonelli free zone / Fubini conditional / failure) (H30). \\
R5 & Parameter measurization & Selmer group, Tate--Shafarevich group
as computable obstructions (H46); Brauer--Manin obstruction. \\
R6 & Parameter programmatization & Cut-elimination as algorithm
(H22 side); Knuth--Bendix completion (H31 meta-layer emergence). \\
\bottomrule
\end{tabular}
\caption{Six recovery patterns with corpus anchors. R1/R2 rows carry
fewer adjudicated exemplars than R3--R6 in the current files; the
mapping is stated at the granularity the evidence supports (per
adjudication discipline, no inflation).}
\end{table}

\subsection{Proposition: Reset, Not Restoration}\label{sec:act2-proposition}

\begin{proposition}[Reset-not-restoration; informal wording, final at S5 fill]
In no corpus instance is the recovery product a copy of the removed
law; in every instance the death position of the old law is occupied
by a structure that keeps its books differently (a new concept
re-wired to the disassembled parameters). Typological induction over
all 29 removal groups.
\end{proposition}

\noindent\emph{Evidence anchor.} The 29-group summary table (Act III
data pack, Table~\ref{tab:act3-summary}) lists the emerged concept
for every group; none restorates the removed law; each realizes one
of R1--R6.
